\subsection{Inducing the shift on three datasets}\label{sec:kinduced}
The contrast between the two official splits is consistent with the shift explanation but does not test
it, because the datasets differ in more than their splits. We therefore induce shift: whole attack
families are removed from the 2{,}000-row training sample and left in the 20{,}000-row test sample, and
the protocol of Table~\ref{tab:twosplits} is repeated (Section~\ref{sec:kdiag}). On UNSW-NB15, ToN-IoT
and CICIDS2017 each family is held out in turn (20 runs), and on UNSW-NB15 two multi-family holdouts are
repeated over three sampling seeds: holdout A removes Exploits, Fuzzers and Reconnaissance ($25\%$ of the
test rows), holdout B Generic, DoS, Analysis and Backdoor ($29\%$). Table~\ref{tab:induced} summarises
the 26 runs and Figure~\ref{fig:induced} shows each of them.

Shift is sufficient to recreate the failure, and its size tracks how badly. Without a holdout, standard
CV tracks test performance on all three datasets (the smallest $r$ over the three kernels is $+0.82$ to
$+0.99$) and loses at most $0.019$ ROC-AUC against the best setting in the grid. With families held out,
the smallest $r$ falls to $+0.46$ on ToN-IoT, $+0.48$ on UNSW-NB15 and $-0.07$ on CICIDS2017, and
selection losses reach $0.061$ for single UNSW-NB15 families, $0.112$ under holdout B and $0.375$ on
CICIDS2017. The label-free squared MMD between the training sample and the unlabelled test sample orders
the runs: it correlates with the smallest $r$ at Spearman $\rho=-0.61$ (permutation $p=0.001$) and with
the largest loss at $\rho=+0.45$ ($p=0.020$); its class-conditional version for the attack rows, which
needs labels, reaches $\rho=-0.77$ and $+0.66$ ($p<0.001$). It is a coarse signal rather than a test:
CICIDS2017 with DDoS held out fails ($r=+0.03$) at a squared MMD of $0.006$, while the holdout-A runs pass
($r\ge+0.83$) at $0.029$ to $0.038$.

What the failure does to a quantum-versus-classical comparison is what matters here. In 5 of the 26 runs
selection reverses the sign of the quantum-minus-best-classical ROC-AUC gap. The largest case recreates
the NSL-KDD artefact on a second dataset: with PortScan held out of CICIDS2017's training sample, the
CV-selected quantum kernel leads the best CV-selected classical kernel by $0.174$ ROC-AUC, while at each
kernel's best setting in the grid it trails by $0.028$. Across all 26 runs the quantum kernel never leads
by more than $0.003$ at the oracle settings. Nor is the outcome stable across training samples: under
holdout B the kernel that selection penalises most changes with the sampling seed (the RBF kernel loses
$0.108$, $0.023$ and $0.101$ over the three seeds, the phase kernel $0.016$, $0.084$ and $0.112$), and
the quantum kernel's lead at $1\%$ false positives on the first seed ($0.301$ against $0.263$ for the RBF
kernel) does not recur on the other two ($0.440$ against $0.664$, and $0.425$ against $0.508$).

Shift-aware selection does not repair the induced failures. Under holdout B (first seed),
leave-attack-types-out CV, with a single usable fold, correlates with the RBF kernel's test ROC-AUC at $r=+0.43$ and chooses a worse
setting than standard CV ($0.751$ against $0.794$): its folds can only hold out attack types the training
sample contains, and a shift towards those is not the shift towards Generic and DoS that the test set
presents. On NSL-KDD the seventeen novel types are variants within the same four attack families as the
seen ones, which is why holding out types imitates the official split there. The result is not an
artefact of the fold construction either. GroupKFold, used throughout, leaves three usable folds on
NSL-KDD and one on UNSW-NB15; a construction that makes all five folds usable gives $r=+0.54$ instead of
$+0.79$ for the RBF kernel on NSL-KDD and $r=-0.38$ under holdout B, and on the unshifted UNSW-NB15 split
it selects an over-smooth RBF setting (test ROC-AUC $0.859$ against $0.938$ under standard CV).
Shift-aware selection is a choice of target rather than a free improvement: it buys robustness to the
kinds of novelty the training data can imitate, at a cost when there is no novelty, and it cannot
anticipate novelty of a kind the training data has never seen. A kernel comparison under distribution
shift should therefore be reported under more than one selection rule, with the oracle bound alongside
and a label-free shift statistic such as the MMD, so that a reader can see how much of a claimed
difference is selection.

@@TAB_INDUCED@@

@@FIG_INDUCED@@
